% !TEX encoding = UTF-8
% !TEX program = xelatex

% 中国科学技术大学专业学位研究生学位申请成果开题报告
% 请先填写 content.tex，再用 XeLaTeX 编译本文件。

\documentclass[UTF8,12pt,fontset=none]{ctexart}

% 页面设置与宏包
\usepackage[paperwidth=184.03mm,paperheight=260.03mm,margin=0mm]{geometry}
\usepackage{tikz,amssymb,framed,fancyhdr}
\usetikzlibrary{calc}

% 优先使用编译平台已安装的字体；平台没有时改用项目 fonts 目录中的同名字体文件。
% 上传到 latex.ustc 时 fonts 目录只需保留 STXINGKA.ttf。
\IfFontExistsTF{Times New Roman}
  {\setmainfont{Times New Roman}}
  {\setmainfont[Path=fonts/,BoldFont=timesbd.ttf,ItalicFont=timesi.ttf,BoldItalicFont=timesbi.ttf]{times.ttf}}
\IfFontExistsTF{SimSun}
  {\setCJKmainfont[BoldFont=SimSun,BoldFeatures={FakeBold=2}]{SimSun}}
  {\setCJKmainfont[Path=fonts/,BoldFont=Simsun.ttc,BoldFeatures={FakeBold=2}]{Simsun.ttc}}
\IfFontExistsTF{SimHei}
  {\setCJKsansfont[AutoFakeBold=2]{SimHei}}
  {\setCJKsansfont[Path=fonts/,AutoFakeBold=2]{SimHei.ttf}}
\IfFontExistsTF{KaiTi}
  {\setCJKfamilyfont{kai}[BoldFont=KaiTi,BoldFeatures={FakeBold=2}]{KaiTi}}
  {\setCJKfamilyfont{kai}[Path=fonts/,BoldFont=Kaiti.ttf,BoldFeatures={FakeBold=2}]{Kaiti.ttf}}
\IfFontExistsTF{STXingkai}
  {\setCJKfamilyfont{xing}[AutoFakeBold=2]{STXingkai}}
  {\setCJKfamilyfont{xing}[Path=fonts/,AutoFakeBold=2]{STXINGKA.ttf}}
\newcommand{\xingkaiword}{\CJKfamily{xing}}
% 封面指定文字加宽汉字间距，不改变字形比例。
\newcommand{\tightcover}{\renewcommand{\CJKglue}{\hskip 0.04em\relax}}
\newcommand{\kaishuword}{\CJKfamily{kai}}
% 说明正文使用楷体，与封面底部使用相同的字体。
\newcommand{\instructionfont}{\CJKfamily{kai}}

% 基本页面样式
\pagestyle{empty}
\setlength{\parindent}{0pt}

% 基本信息和报告正文请在 content.tex 中填写。
% !TEX root = main.tex
% 在此填写开题报告的信息与正文；编译时使用 main.tex。
% 各命令中预填的是表格原有的提示文本，填写时请删除后再写入自己的内容。

% ==================== 基本信息 ====================
% 简况表和封面为固定布局，填写较长内容后请检查是否超出边框。
% 封面题目超出第一条横线时会自动续写到第二条横线；两行仍放不下时自动缩小题目字号。
\newcommand{\reporttitle}{}
\newcommand{\studentname}{}
\newcommand{\studentid}{}
\newcommand{\gender}{}
\newcommand{\supervisor}{}
\newcommand{\secondsupervisor}{（选填）}
\newcommand{\practicesupervisor}{}
\newcommand{\department}{软件学院}
\newcommand{\majorcategory}{电子信息 0854}
\newcommand{\majorfield}{软件工程 085405}
\newcommand{\reportdate}{}

% ==================== 勾选项 ====================
% 填写选项名称即可勾选，名称须与下方列出的完全一致，多个选项用英文逗号分隔；留空则不勾选。
% 培养方式：全日制、非全日制
\newcommand{\studymode}{}
% 专业学位类型：硕士、博士
\newcommand{\degreetype}{}
% 硕士成果形式：专题研究类论文、调研报告、案例分析报告、产品设计（作品创作）、方案设计、其他
\newcommand{\masterachievement}{}
\newcommand{\masterother}{}% 勾选“其他”时填写具体形式
% 博士成果形式：专题研究类论文、重大装备、仪器设备、其他硬件产品、软件产品、设计方案、技术标准、其他
\newcommand{\doctorachievement}{}
\newcommand{\doctorother}{}% 勾选“其他”时填写具体形式

\newcommand{\englishtitle}{}
\newcommand{\abstracttext}{}
% 中文主题词之间请用 \quad 分隔，如 {异常检测\quad 智能运维}；直接输入的空格会被忽略。
\newcommand{\keywords}{}
\newcommand{\englishkeywords}{}
\newcommand{\sourceoftopic}{%
  % 提示文本，填写时删除。
  工程博士必须填写研究所依托的重大工程项目，并标明项目编号。
}

% ==================== 报告正文 ====================
% 在各命令的花括号内填写内容，段落之间可用 \par。
% 第二至第七节的边框可随内容自动跨页。

% 二、选题依据
\newcommand{\topicbasis}{%
  % 提示文本，填写时删除。
  阐述该选题的研究意义，或工程设计的价值和意义，国内外概况和发展趋势，选题的先进性和实用性，技术难度及工作量。附上主要参考文献（列出作者、论文名称、期刊名称、出版年月）。
}

% 三、研究的创新之处
\newcommand{\innovation}{%
  % 提示文本，填写时删除。
  研究内容、拟采用的研究方法、技术路线等方面有哪些创新之处（博士必填，硕士选填）。
}

% 四、研究方法与技术路线
\newcommand{\methodroute}{%
  % 提示文本，填写时删除。
  拟采用的设计方法，技术路线，实践方案的可行性分析。
}

% 五、预期成果
\newcommand{\expectedresults}{%
  % 提示文本，填写时删除。
  预期达到的目标和取得的成果。
}

% 六、已取得的与学位申请成果研究内容相关的成果
\newcommand{\existingresults}{%
  % 提示文本，填写时删除。
  已发表或被接收发表的学术论文或其它相关创新性成果。
}

% 七、研究计划与进度安排
\newcommand{\researchplan}{%
  % 提示文本，填写时删除。
  应包括文献调研，工程设计，新工艺、新材料、新设备、新产品的研制和调试，实验操作，实验数据的分析处理，撰写论文等。
}


% ==================== 页面排版命令 ====================
% 以下命令控制固定页面的坐标和边框，通常无需修改。
% 勾选框优先使用 DejaVu Sans 的 ☐ / ☑；没有该字体时用 □ 加 TikZ 绘制的勾。
% 两种方案都按原 □ 的尺寸占位，勾选与否、切换字体都不影响表格排版。
% 方框下移 0.075em，使其与汉字的视觉中心对齐（方框底边原本落在基线上，比汉字偏高）。
% 为避免 xeCJK 在勾选框前后插入中西文间距：外层用 \hbox 封装，DejaVu 字形用 \XeTeXglyph 直接放置。
\newsavebox{\boxglyphref}
\newsavebox{\checkglyphref}
\newcommand{\boxglyph}[1]{\sbox{\boxglyphref}{\scalebox{1.25}{$\square$}}%
  \hbox{\makebox[\wd\boxglyphref][c]{\raisebox{-.075em}[\ht\boxglyphref][\dp\boxglyphref]{#1}}}}
\IfFontExistsTF{DejaVu Sans}
  {\newfontfamily\checkboxfont{DejaVu Sans}
   \newcommand{\boxmark}{\boxglyph{\scalebox{1.2}{\checkboxfont\mdseries\XeTeXglyph\XeTeXcharglyph"2610\relax}}}
   \newcommand{\checkedbox}{\boxglyph{\scalebox{1.2}{\checkboxfont\mdseries\XeTeXglyph\XeTeXcharglyph"2611\relax}}}}
  {\newcommand{\boxmark}{\boxglyph{\scalebox{1.25}{$\square$}}}
   \newcommand{\checkedbox}{\boxglyph{\scalebox{1.25}{\sbox{\checkglyphref}{$\square$}%
     \rlap{\usebox{\checkglyphref}}\makebox[\wd\checkglyphref][c]{\tikz[baseline=0pt,x=\wd\checkglyphref,y=\wd\checkglyphref]\draw[line width=.45pt,line cap=round,line join=round] (.22,.42)--(.42,.2)--(.8,.78);}}}}}
% \choice{已选列表}{选项}：选项在 content.tex 的列表中时显示 ☑，否则显示 □。
% \checkchoices{已选列表}{可选项}：列表中有无法识别的选项时给出警告。
% \otherblank{内容}：“其他”后的横线，内容较长时横线随文字加长。
\ExplSyntaxOn
\clist_new:N \l__kt_choice_clist
\msg_new:nnn { kt } { unknown-choice } { 无法识别的勾选项“#1”，可选项为：#2 }
\NewDocumentCommand{\choice}{mm}
  {
    \clist_set:No \l__kt_choice_clist {#1}
    \clist_if_in:NnTF \l__kt_choice_clist {#2} { \checkedbox } { \boxmark }
  }
\NewDocumentCommand{\checkchoices}{mm}
  {
    \clist_set:No \l__kt_choice_clist {#1}
    \clist_map_inline:Nn \l__kt_choice_clist
      { \clist_if_in:nnF {#2} {##1} { \msg_warning:nnnn { kt } { unknown-choice } {##1} {#2} } }
  }
\NewDocumentCommand{\otherblank}{m}
  {
    \hbox_set:Nn \l_tmpa_box {#1}
    \dim_compare:nNnTF { \box_wd:N \l_tmpa_box } > { 18mm }
      { \underline{\hspace{1mm}#1\hspace{1mm}} }
      { \underline{\makebox[20mm]{#1}} }
  }
\ExplSyntaxOff
\newcommand{\pagebegin}{\begin{tikzpicture}[remember picture,overlay,x=1mm,y=-1mm,shift={(current page.north west)},every node/.style={inner sep=0pt,outer sep=0pt}]}
\newcommand{\pageend}{\end{tikzpicture}\null\newpage}
\newcommand{\txt}[4]{\node[anchor=north west,text width=#3mm,align=left] at (#1,#2) {#4};}
\newcommand{\centertext}[4]{\node[anchor=north,text width=#3mm,align=center] at (#1,#2) {#4};}
\newcommand{\cellcenter}[4]{\node[anchor=center,text width=#3mm,align=center] at (#1,#2) {\fontsize{10}{14}\selectfont #4\par};}
\newcommand{\cellleft}[4]{\node[anchor=west,text width=#3mm,align=left] at (#1,#2) {\fontsize{10}{14}\selectfont #4\par};}
\newcommand{\border}[4]{\draw[line width=.3pt] (#1,#2) rectangle ++(#3,#4);}
\newcommand{\heading}[3]{\txt{#1}{#2}{155}{\sffamily\fontsize{15}{18}\selectfont #3}}
\newcommand{\inside}[4]{\txt{#1}{#2}{#3}{\fontsize{10}{14}\selectfont #4}}
\newcommand{\num}[1]{\centertext{92}{246}{10}{\fontsize{7}{8}\selectfont #1}}
% 封面题目：一行放得下时居中于第一条横线；放不下时自动续写到第二条横线。
% 两行仍放不下时，题目字号从 22pt 起每次缩小 0.5pt，直到恰好放进两行（最小 14pt）。
\newsavebox{\covertitlebox}
\newlength{\covertitlewidth}
\newlength{\covertitlesize}
\newcount\covertitlelines
\newcommand{\covertitlelabel}{学位申请成果题目：}
% 封面不显示“（样表）”时，题目与信息栏整体上移 7mm，使空出的半行留白上下均分。
\newsavebox{\coversamplebox}
\newcommand{\covershift}{0}
\newcommand{\covertitlefont}{\fontsize{\the\covertitlesize}{\the\covertitlesize}\selectfont}
% 按封面上的实际排法试排一次，记下题目占用的行数。
\newcommand{\covertitlecount}{%
  \setbox0\vbox{\hsize=\covertitlewidth\parindent=0pt\raggedright
    \fontsize{22}{26}\selectfont\setlength{\baselineskip}{11.6mm}%
    \covertitlelabel{\covertitlefont\reporttitle}\par\global\covertitlelines=\prevgraf}}
\newcommand{\covertitleshrink}{%
  \covertitlecount
  \let\next\relax
  \ifnum\covertitlelines>2 \ifdim\covertitlesize>14pt
    \addtolength{\covertitlesize}{-0.5pt}\let\next\covertitleshrink
  \fi\fi
  \next}
\begin{document}

% ==================== 封面 ====================
\sbox{\covertitlebox}{\fontsize{22}{26}\selectfont\reporttitle}
\settowidth{\covertitlewidth}{\fontsize{22}{26}\selectfont\covertitlelabel}
\addtolength{\covertitlewidth}{65mm}
\setlength{\covertitlesize}{22pt}
\sbox{\coversamplebox}{\coversample}
\ifdim\wd\coversamplebox=0pt \renewcommand{\covershift}{-7}\fi
\ifdim\wd\covertitlebox>65mm
  \covertitleshrink
  \ifnum\covertitlelines>2
    \PackageWarningNoLine{kt}{封面题目在 14pt 下仍超过两行，请缩短题目}
  \else\ifdim\covertitlesize<22pt
    \typeout{封面题目较长，字号已缩小为 \the\covertitlesize}
  \fi\fi
\fi
\pagebegin
\centertext{92}{18}{165}{\xingkaiword\fontsize{26}{30}\selectfont \tightcover 中国科学技术大学}
\centertext{92}{33}{170}{\sffamily\fontsize{24}{28}\selectfont 专业学位研究生学位申请成果开题报告}
\centertext{92}{47}{100}{\fontsize{26}{30}\selectfont\coversample}
\ifdim\wd\covertitlebox>65mm
  \node[anchor=north west] (cover-title-line) at (30,61+\covershift) {\fontsize{22}{26}\selectfont \phantom{\covertitlelabel}\underline{\makebox[65mm]{}}};
  \node[anchor=north west,text width=\covertitlewidth,align=left] at (30,61+\covershift) {\fontsize{22}{26}\selectfont\setlength{\baselineskip}{11.6mm}\covertitlelabel{\covertitlefont\reporttitle}\par};
\else
  \node[anchor=north west] (cover-title-line) at (30,61+\covershift) {\fontsize{22}{26}\selectfont \covertitlelabel\underline{\makebox[65mm][c]{\reporttitle}}};
\fi
\draw[line width=.55pt] (30,80+\covershift)--(cover-title-line.east |- 0,80+\covershift);
\foreach \y/\lab/\val in {87/学生姓名/\studentname,101/学生学号/\studentid,115/校内导师/\supervisor,129/第二导师/\secondsupervisor,143/实践导师/\practicesupervisor,157/所在院系/\department,171/专业类别/\majorcategory,185/专业领域/\majorfield,199/填表日期/\reportdate}{\txt{30}{\y+\covershift}{140}{\fontsize{22}{26}\selectfont \lab：\underline{\makebox[95mm][c]{\val}}}}
\centertext{92}{219}{160}{\kaishuword\bfseries\fontsize{15}{18}\selectfont \tightcover 中国科学技术大学研究生院培养处}
\centertext{92}{231}{120}{\kaishuword\bfseries\fontsize{15}{18}\selectfont \tightcover 二零二五年三月\quad 制表}
\pageend
% ==================== 说明 ====================
\pagebegin
\centertext{92}{28}{70}{\rmfamily\bfseries\fontsize{22}{26}\selectfont 说\quad 明}
\txt{15}{50}{154}{\instructionfont\fontsize{18}{24}\selectfont \hspace*{12mm}1.\quad 学位申请成果的开题报告是保证研究生培养质量的一个重要环节，为了加强对研究生培养的过程管理，规范其学位申请成果的开题报告，特制此表。\par}
\txt{15}{83}{154}{\instructionfont\fontsize{18}{24}\selectfont \hspace*{12mm}2.\quad 学位申请成果开题报告，应该在学位授予点或培养单位组织的学术报告会上报告，听取意见，论证后再填写此表。\par}
\txt{15}{116}{154}{\instructionfont\fontsize{18}{24}\selectfont \hspace*{12mm}3.\quad 此表经导师签字后，交培养单位研究生教学管理办公室存档。\par}
\txt{15}{138}{154}{\instructionfont\fontsize{18}{24}\selectfont \hspace*{12mm}4.\quad 研究生在提交学位申请时，必须提交该学位申请成果开题报告。\par}
\pageend
% ==================== 一、简况 ====================
% 外框与合并单元格按 Word 参照页定位。
\pagebegin
\heading{15}{17}{一、简况}
% 表格外框和分割线
\border{13.5}{27}{157}{214}
\foreach \y in {56,81,108,131,214} {\draw[line width=.3pt] (13.5,\y)--(170.5,\y);}
\draw[line width=.3pt] (20.3,41)--(170.5,41) (20.3,189)--(170.5,189);
\draw[line width=.3pt] (20.3,27)--(20.3,56) (39.8,27)--(39.8,56) (69.8,27)--(69.8,41) (94.7,27)--(94.7,41) (134,27)--(134,41) (151,27)--(151,41);
\draw[line width=.3pt] (30,56)--(30,131) (53.5,108)--(53.5,131) (30,119.5)--(170.5,119.5);
\draw[line width=.3pt] (20.3,131)--(20.3,241) (30,131)--(30,214) (30,200)--(170.5,200) (43.5,200)--(43.5,214) (30,207)--(170.5,207);

% 研究生信息与培养方式
\cellcenter{16.9}{41.5}{5}{\sffamily\bfseries\fontsize{10}{12}\selectfont 研\\究\\生\\简\\况}
\cellcenter{30}{34}{18}{\fontsize{10}{12}\selectfont 姓\quad 名}
\cellcenter{55}{34}{28}{\fontsize{10}{12}\selectfont \studentname}
\cellcenter{82}{34}{20}{\fontsize{10}{12}\selectfont 学\quad 号}
\cellcenter{114}{34}{35}{\fontsize{10}{12}\selectfont \studentid}
\cellcenter{142}{34}{16}{\fontsize{10}{12}\selectfont 性\quad 别}
\cellcenter{160}{34}{17}{\fontsize{10}{12}\selectfont \gender}
\cellcenter{30}{48.5}{18}{\fontsize{10}{12}\selectfont 培养方式}
\checkchoices{\studymode}{全日制,非全日制}\checkchoices{\degreetype}{硕士,博士}
\checkchoices{\masterachievement}{专题研究类论文,调研报告,案例分析报告,产品设计（作品创作）,方案设计,其他}
\checkchoices{\doctorachievement}{专题研究类论文,重大装备,仪器设备,其他硬件产品,软件产品,设计方案,技术标准,其他}
\cellleft{62}{48.5}{100}{\choice{\studymode}{全日制}\ 全日制\hspace{17mm}\choice{\studymode}{非全日制}\ 非全日制}

% 专业学位类型
\cellcenter{21.75}{68.5}{15}{\sffamily\bfseries\fontsize{10}{12}\selectfont \choice{\degreetype}{硕士}\\专业学位\\硕士}
\cellleft{34}{68.5}{133}{\makebox[33mm][l]{\choice{\masterachievement}{专题研究类论文}\ 专题研究类论文}\makebox[26mm][l]{\choice{\masterachievement}{调研报告}\ 调研报告}\makebox[29mm][l]{\choice{\masterachievement}{案例分析报告}\ 案例分析报告}\choice{\masterachievement}{产品设计（作品创作）}\ 产品设计（作品创作）\\[2mm] \makebox[33mm][l]{\choice{\masterachievement}{方案设计}\ 方案设计}\choice{\masterachievement}{其他}\ 其他\ \otherblank{\masterother}}
\cellcenter{21.75}{94.5}{15}{\sffamily\bfseries\fontsize{10}{12}\selectfont \choice{\degreetype}{博士}\\专业学位\\博士}
\cellleft{34}{94.5}{133}{学位论文申请：\choice{\doctorachievement}{专题研究类论文}\ 专题研究类论文\\[2mm] \makebox[27mm][l]{实践成果申请：}\makebox[25mm][l]{\choice{\doctorachievement}{重大装备}\ 重大装备}\makebox[26mm][l]{\choice{\doctorachievement}{仪器设备}\ 仪器设备}\makebox[35mm][l]{\choice{\doctorachievement}{其他硬件产品}\ 其他硬件产品}\choice{\doctorachievement}{软件产品}\ 软件产品\\[2mm] \makebox[27mm][l]{}\makebox[25mm][l]{\choice{\doctorachievement}{设计方案}\ 设计方案}\makebox[26mm][l]{\choice{\doctorachievement}{技术标准}\ 技术标准}\choice{\doctorachievement}{其他}\ 其他\ \otherblank{\doctorother}}

% 成果题目、摘要和关键词
\cellcenter{21.75}{119.5}{15}{\sffamily\bfseries\fontsize{10}{12}\selectfont 学位申请\\成果题目}
\cellcenter{41.75}{113.75}{20}{\fontsize{10}{12}\selectfont 中\quad 文}
\cellcenter{41.75}{125.25}{20}{\fontsize{10}{12}\selectfont 英\quad 文}
\cellleft{55}{113.75}{110}{\reporttitle}\cellleft{55}{125.25}{110}{\englishtitle}
\cellcenter{16.9}{172.5}{5}{\sffamily\bfseries\fontsize{10}{10}\selectfont 学\\位\\申\\请\\成\\果\\摘\\要\\和\\关\\键\\词}
\cellcenter{25}{160}{8}{\fontsize{10}{14}\selectfont 摘\\要}
\inside{32}{133}{135}{\abstracttext}
\cellcenter{25}{201.5}{8}{\fontsize{10}{12}\selectfont 关\\键\\词}
\cellcenter{36.75}{203.5}{10}{中文}\cellcenter{36.75}{210.5}{10}{英文}
\cellleft{31}{194.5}{138}{关键词数量三至五个，主题词之间空一格（英文用“/ ”分隔）}
\cellleft{45}{203.5}{120}{\keywords}\cellleft{45}{210.5}{120}{\englishkeywords}

% 选题来源
\cellcenter{17}{227.5}{5}{\sffamily\bfseries\fontsize{10}{12}\selectfont 选\\题\\来\\源}
\inside{21}{215}{145}{\sourceoftopic}
\num{1}\pageend
% ==================== 二至七、报告正文 ====================
% 正文以可分页边框排版；宏内可使用 \par 分段。
% 空白时保留与原表近似的留白；填入内容后边框随文字跨页。
\newgeometry{left=16mm,right=16mm,top=17mm,bottom=20mm}
\pagestyle{fancy}\fancyhf{}\fancyfoot[C]{\fontsize{9}{11}\selectfont\thepage}\renewcommand{\headrulewidth}{0pt}
\setcounter{page}{2}
\setlength{\parindent}{0pt}
\newcommand{\flowheading}[1]{\par\noindent{\sffamily\fontsize{15}{19}\selectfont #1}\par\vspace{1mm}}
\newlength{\flowremaining}
\newsavebox{\flowmeasure}
\newcommand{\flowbox}[2]{%
  \setlength{\FrameRule}{0.3pt}\setlength{\FrameSep}{2mm}\setlength{\topsep}{0pt}\setlength{\partopsep}{0pt}%
  % 以空白模板的留白为最小高度；正文超出时由 framed 自动分页。
  \setlength{\flowremaining}{#2}%
  \sbox{\flowmeasure}{\begin{minipage}{\dimexpr\linewidth-2\FrameSep-2\FrameRule\relax}%
    \fontsize{10.5}{15}\selectfont #1\par\end{minipage}}%
  \addtolength{\flowremaining}{-\ht\flowmeasure}%
  \addtolength{\flowremaining}{-\dp\flowmeasure}%
  \begin{framed}\fontsize{10.5}{15}\selectfont #1\par
  \ifdim\flowremaining>0pt\vspace*{\flowremaining}\fi
  \end{framed}\vspace{1mm}}

% 二、选题依据
\flowheading{二、选题依据}
\flowbox{\topicbasis}{129.12mm}

% 三、研究的创新之处
\flowheading{三、研究的创新之处}
\flowbox{\innovation}{55.43mm}
\clearpage

% 四、研究方法与技术路线
\flowheading{四、研究方法与技术路线}
\flowbox{\methodroute}{129.43mm}

% 五、预期成果
\flowheading{五、预期成果}
\flowbox{\expectedresults}{51.42mm}
\clearpage

% 六、已取得的与学位申请成果研究内容相关的成果
\flowheading{六、已取得的与学位申请成果研究内容相关的成果}
\flowbox{\existingresults}{60.42mm}
\vspace{4mm}

% 七、研究计划与进度安排
\flowheading{七、研究计划与进度安排}
\flowbox{\researchplan}{68.27mm}

% 签名
\vspace{12mm}
\hfill\begin{minipage}{90mm}\fontsize{12}{22}\selectfont
研究生本人签名：\hrulefill\\
校内导师（主导师）签名：\hrulefill\\
实践导师签名：\hrulefill\\
\makebox[\linewidth][r]{年\qquad 月\qquad 日}
\end{minipage}
\end{document}
